\documentclass[12pt, a4paper]{article}

% ── Codificación y lenguaje ───────────────────────────────────────────────────
\usepackage[utf8]{inputenc}
\usepackage[T1]{fontenc}
\usepackage[spanish]{babel}

% ── Márgenes ─────────────────────────────────────────────────────────────────
\usepackage[left=2.5cm, right=2.5cm, top=2.5cm, bottom=2.5cm]{geometry}

% ── Matemáticas ──────────────────────────────────────────────────────────────
\usepackage{amsmath, amssymb}

% ── Código fuente ────────────────────────────────────────────────────────────
\usepackage{listings}
\usepackage{xcolor}

\definecolor{codebg}{RGB}{245,245,245}
\definecolor{codeframe}{RGB}{200,200,200}
\definecolor{keywordcolor}{RGB}{0,0,180}
\definecolor{commentcolor}{RGB}{60,130,60}
\definecolor{stringcolor}{RGB}{180,0,0}

\lstdefinestyle{cpp}{
  language=C++,
  backgroundcolor=\color{codebg},
  frame=single,
  rulecolor=\color{codeframe},
  basicstyle=\ttfamily\small,
  keywordstyle=\color{keywordcolor}\bfseries,
  commentstyle=\color{commentcolor}\itshape,
  stringstyle=\color{stringcolor},
  numbers=left,
  numberstyle=\tiny\color{gray},
  numbersep=8pt,
  breaklines=true,
  breakatwhitespace=false,
  tabsize=4,
  showstringspaces=false,
  captionpos=b,
  literate=
    {á}{{\'a}}1 {é}{{\'e}}1 {í}{{\'i}}1 {ó}{{\'o}}1 {ú}{{\'u}}1
    {Á}{{\'A}}1 {É}{{\'E}}1 {Í}{{\'I}}1 {Ó}{{\'O}}1 {Ú}{{\'U}}1
    {ñ}{{\~n}}1 {Ñ}{{\~N}}1
    {¿}{{?`}}1 {¡}{{!`}}1,
}

\lstdefinestyle{pseudo}{
  backgroundcolor=\color{codebg},
  frame=single,
  rulecolor=\color{codeframe},
  basicstyle=\ttfamily\small,
  numbers=left,
  numberstyle=\tiny\color{gray},
  numbersep=8pt,
  breaklines=true,
  tabsize=4,
  showstringspaces=false,
  captionpos=b,
  literate=
    {á}{{\'a}}1 {é}{{\'e}}1 {í}{{\'i}}1 {ó}{{\'o}}1 {ú}{{\'u}}1
    {Á}{{\'A}}1 {É}{{\'E}}1 {Í}{{\'I}}1 {Ó}{{\'O}}1 {Ú}{{\'U}}1
    {ñ}{{\~n}}1 {Ñ}{{\~N}}1,
}

% ── Tablas ───────────────────────────────────────────────────────────────────
\usepackage{booktabs}
\usepackage{array}
\usepackage{longtable}

% ── Gráficos / colores ───────────────────────────────────────────────────────
\usepackage{graphicx}

% ── Hipervínculos ────────────────────────────────────────────────────────────
\usepackage{hyperref}
\hypersetup{
  colorlinks=true,
  linkcolor=blue,
  urlcolor=blue,
  citecolor=blue,
}

% ── Encabezado y pie de página ────────────────────────────────────────────────
\usepackage{fancyhdr}
\pagestyle{fancy}
\fancyhf{}
\rhead{Trabajo Práctico 01}
\lhead{Universidad EAFIT --- Análisis y Diseño de Algoritmos}
\cfoot{\thepage}

% ── Título ───────────────────────────────────────────────────────────────────
\title{
  \textbf{ESCUELA DE CIENCIAS APLICADAS E INGENIERÍA}\\
  \large Departamento de Ingeniería\\[0.5em]
  \Large Trabajo Práctico 01\\[0.3em]
  \LARGE Ejercicio 1: Permutaciones con Restricciones (Fuerza Bruta)\\[0.5em]
  \normalsize Asignatura: Análisis y Diseño de Algoritmos\\
  Profesor: Carlos Alberto Álvarez Henao
}
\author{
  Andrés Felipe Vélez Álvarez \\
  Sebastián Salazar Henao \\
  Nathalia Valentina Cardoza Azuaje
}
\date{5 de abril de 2026 \\ Universidad EAFIT}

% ─────────────────────────────────────────────────────────────────────────────
\begin{document}

\maketitle
\thispagestyle{empty}
\newpage

\tableofcontents
\newpage

% ══════════════════════════════════════════════════════════════════════════════
\section{Descripción Teórica del Problema}
% ══════════════════════════════════════════════════════════════════════════════

Dado un conjunto de $n$ enteros positivos distintos $A = \{a_0, a_1, \ldots, a_{n-1}\}$,
se desea generar todas las permutaciones posibles del conjunto y determinar cuáles
satisfacen la siguiente restricción de adyacencia:

\begin{equation}
  P[i] \leq 2 \cdot P[i+1] \quad \forall\, i \in \{0, 1, \ldots, n-2\}
  \label{eq:restriccion}
\end{equation}

Es decir, en toda permutación válida ningún elemento puede ser mayor que el doble
del elemento que le sigue. El programa debe: (a)~generar todas las permutaciones
del conjunto mediante fuerza bruta (enumeración exhaustiva del factorial de
permutaciones), (b)~filtrar las que cumplen la restricción, y (c)~reportar el total
de permutaciones generadas, el total de válidas y el listado de éstas.

El número total de permutaciones de $n$ elementos distintos es $n!$, lo que determina
directamente la complejidad del enfoque. A modo de referencia:

\begin{center}
\begin{tabular}{rrl}
\toprule
$n$ & $n!$ & Tiempo estimado \\
\midrule
5  & 120            & Instantáneo    \\
8  & 40\,320        & Muy rápido     \\
10 & 3\,628\,800    & Rápido         \\
12 & 479\,001\,600  & Lento          \\
13 & 6\,227\,020\,800 & Impracticable \\
\bottomrule
\end{tabular}
\end{center}

% ══════════════════════════════════════════════════════════════════════════════
\section{Enfoques de Solución}
% ══════════════════════════════════════════════════════════════════════════════

\subsection{Fuerza Bruta}

La fuerza bruta aborda el problema mediante enumeración exhaustiva: se generan
absolutamente todas las $n!$ permutaciones del conjunto y, una vez construida cada
una, se verifica si cumple la restricción $P[i] \leq 2 \cdot P[i+1]$ para todos sus pares
consecutivos. Las permutaciones que pasan el filtro se cuentan e imprimen; las que
no, se descartan.

\textbf{Proceso paso a paso:}
\begin{enumerate}
  \item Se lee el conjunto de $n$ enteros y se ordena de forma ascendente. Esto
        garantiza comenzar desde la permutación lexicográficamente menor, condición
        necesaria para que \texttt{next\_permutation} recorra exactamente las $n!$ permutaciones.
  \item Se itera sobre todas las permutaciones mediante un bucle \texttt{do-while}
        combinado con la función \texttt{next\_permutation}.
  \item En cada iteración se verifica la restricción recorriendo los $n-1$ pares
        consecutivos. Si algún par la viola, la permutación se descarta con \texttt{break}.
  \item Al terminar, se reportan los totales y el listado de permutaciones válidas.
\end{enumerate}

\subsection{Por qué este enfoque es Fuerza Bruta y no Backtracking}

Aunque ambos algoritmos resuelven el mismo problema y producen los mismos
resultados, la diferencia fundamental radica en \textit{cuándo} se evalúa la restricción
durante la construcción de cada permutación.

En la fuerza bruta implementada, la permutación se construye \textbf{completamente}
antes de verificar si es válida. Es decir, se genera la permutación entera de $n$
elementos y sólo al final se recorren sus pares consecutivos para aceptarla o
descartarla. En ningún momento se interrumpe la construcción anticipadamente: si los
primeros dos elementos ya violan la restricción, el algoritmo termina de armar la
permutación completa antes de detectarlo y descartarla.

En el \textit{backtracking}, en cambio, la verificación ocurre \textbf{durante} la
construcción, posición por posición. Al colocar el elemento en la posición $i$, se
comprueba inmediatamente si viola la restricción con el elemento en $i-1$. Si la
viola, se poda la rama completa: no se exploran ninguna de las $(n-i-1)!$
sub-permutaciones que se derivan de ese prefijo inválido.

La clave definitoria es la \textbf{ausencia de poda}:

\begin{itemize}
  \item \textbf{Fuerza Bruta:} genera las $n!$ permutaciones sin excepción, evalúa la
        restricción \textit{a posteriori} sobre la permutación completa.
  \item \textbf{Backtracking:} construye parcialmente y poda \textit{a priori} las ramas
        que no pueden llevar a soluciones válidas.
\end{itemize}

Por tanto, aunque el código utilice \texttt{next\_permutation} (que internamente aplica
un algoritmo de generación lexicográfica eficiente), la estrategia de búsqueda sigue
siendo fuerza bruta: el espacio de búsqueda explorado es siempre el total de las $n!$
permutaciones, sin ninguna reducción anticipada basada en la restricción del problema.

\subsection{Justificación del uso de \texttt{next\_permutation} de la STL}

Para generar las permutaciones se utiliza la función \texttt{std::next\_permutation}
de la biblioteca estándar de C++ (\texttt{<algorithm>}) en lugar de implementar un
algoritmo propio. Esta decisión se justifica por las siguientes razones:

\begin{enumerate}
  \item \textbf{Corrección garantizada.} La implementación de la STL ha sido
        ampliamente probada y verificada. Una implementación propia del algoritmo de
        generación lexicográfica puede introducir errores sutiles, por ejemplo en el
        manejo de los índices de frontera o en la inversión del sufijo.

  \item \textbf{Completitud del recorrido.} \texttt{next\_permutation} garantiza
        recorrer exactamente las $n!$ permutaciones en orden lexicográfico estricto,
        sin repetir ni omitir ninguna, siempre que el arreglo de entrada esté ordenado
        ascendentemente al inicio del bucle.

  \item \textbf{Eficiencia.} El algoritmo interno opera en tiempo amortizado $O(n)$
        por llamada, que es el óptimo posible para generar la siguiente permutación en
        orden lexicográfico.

  \item \textbf{Legibilidad.} El patrón \texttt{do-while} con \texttt{next\_permutation}
        expresa de forma directa la intención del algoritmo sin oscurecer la lógica
        principal con detalles de implementación del generador.

  \item \textbf{Portabilidad.} Al ser parte del estándar ISO de C++, la función está
        disponible en cualquier compilador compatible sin dependencias externas.
\end{enumerate}

La implementación propia \texttt{next\_permutation\_test} se incluye en el código
exclusivamente con fines educativos, para evidenciar el mecanismo interno del
algoritmo de generación lexicográfica.

\subsection{Backtracking (referencia comparativa)}

El backtracking mejoraría el enfoque construyendo la permutación posición por posición
y verificando la restricción en cada paso antes de agregar el siguiente elemento. Si en
la posición $i$ la condición $A[i-1] \leq 2 \cdot A[i]$ no se cumple, se poda esa rama
completa sin explorar ninguna de sus sub-permutaciones. A esta eliminación anticipada
se le llama \textbf{poda}.

\subsection{Comparación de Enfoques}

\begin{center}
\begin{tabular}{lll}
\toprule
\textbf{Aspecto}            & \textbf{Fuerza Bruta}         & \textbf{Backtracking}                  \\
\midrule
Espacio de búsqueda         & Todas las $n!$ permutaciones  & Árbol podado                           \\
Poda                        & No                            & Sí (viola restricción $\Rightarrow$ descarta) \\
Completitud                 & Sí                            & Sí                                     \\
Corrección                  & Sí                            & Sí                                     \\
Complejidad peor caso       & $O(n! \cdot n)$               & $\leq O(n! \cdot n)$                   \\
Eficiencia práctica         & Inferior para $n$ grande      & Mejor con restricciones fuertes        \\
\bottomrule
\end{tabular}
\end{center}

% ══════════════════════════════════════════════════════════════════════════════
\section{Pseudocódigo}
% ══════════════════════════════════════════════════════════════════════════════

\subsection{Fuerza Bruta}

\begin{lstlisting}[style=pseudo, caption={Pseudocódigo: Permutaciones con Restricción por Fuerza Bruta}]
leer n
leer A[0..n-1]   // conjunto de enteros distintos
ordenar A asc    // garantiza comenzar desde la permutación menor

total_generadas = 0
total_validas   = 0

// Calcular n! para verificación
cont = 1
para i en [1..n]:
    cont = cont * i

// Bucle principal: recorre todas las n! permutaciones
hacer:
    total_generadas += 1
    valida = verdadero

    // Verificar la restricción en cada par consecutivo
    para i en [0..n-2]:
        si A[i] > 2 * A[i+1]:
            valida = falso
            break   // no hace falta seguir verificando

    si valida:
        imprimir A
        total_validas += 1

mientras next_permutation(A) exista

imprimir "Total permutaciones posibles: ",  cont
imprimir "Total permutaciones generadas:", total_generadas
imprimir "Total permutaciones válidas: ",  total_validas
\end{lstlisting}

La función \texttt{next\_permutation} recibe el arreglo y lo reorganiza en su siguiente
permutación en orden lexicográfico, modificándolo directamente en memoria. Retorna
\texttt{true} si existe una permutación siguiente, o \texttt{false} cuando ya se alcanzó
la permutación máxima (orden descendente).

\subsection{Algoritmo interno de \texttt{next\_permutation}}

\begin{lstlisting}[style=pseudo, caption={Pseudocódigo interno: next\_permutation}]
next_permutation(inicio, fin):
    n = fin - inicio
    si n <= 1: retornar falso

    // Paso 1: buscar punto de quiebre de derecha a izquierda
    i = n - 2
    mientras i >= 0 y inicio[i] >= inicio[i+1]:
        i = i - 1

    si i == -1: retornar falso   // última permutación (todo descendente)

    // Paso 2: buscar el primer elemento mayor que inicio[i] desde la derecha
    j = n - 1
    mientras inicio[j] <= inicio[i]:
        j = j - 1

    // Paso 3: intercambiar inicio[i] e inicio[j]
    swap(inicio[i], inicio[j])

    // Paso 4: invertir el sufijo desde i+1 hasta el final
    invertir(inicio[i+1..n-1])

    retornar verdadero
\end{lstlisting}

% ══════════════════════════════════════════════════════════════════════════════
\section{Implementación en C++}
% ══════════════════════════════════════════════════════════════════════════════

La implementación completa se encuentra en el archivo \texttt{main\_ejercicio1.cpp}
adjunto al presente informe. El código se organiza en los siguientes bloques:

\begin{itemize}
  \item \textbf{\texttt{main}:} lectura del conjunto, ordenamiento mediante burbuja,
        bucle principal con \texttt{do-while}, medición de tiempo y reporte de resultados.
  \item \textbf{\texttt{factorial}:} calcula $n!$ de forma iterativa para reportar el
        total teórico y verificar que \texttt{total\_generadas == factorial(n)}.
  \item \textbf{\texttt{next\_permutation\_test}:} implementación propia del algoritmo
        de generación lexicográfica, incluida para evidenciar el mecanismo interno.
  \item \textbf{Medición de tiempo:} se usa \texttt{chrono::high\_resolution\_clock}
        para medir únicamente el tiempo del bucle principal.
\end{itemize}

Para compilar y ejecutar:
\begin{lstlisting}[style=pseudo]
g++ -O2 -std=gnu++17 -o solucion permutaciones.cpp
./solucion
\end{lstlisting}

\begin{lstlisting}[style=cpp, caption={Implementación completa: permutaciones.cpp}]
#include <iostream>
#include <algorithm>
#include <chrono>
using namespace std;

int  factorial(int n);
bool next_permutation_test(int* inicio, int* fin);

int main() {

    int n = 0, total_generadas = 0, total_validas = 0, aux = 0;
    bool valida;

    cout << "Ingrese el tamaño del arreglo: ";
    cin >> n;
    int A[n];

    for (int i = 0; i < n; i++) {
        cout << "Ingrese el elemento " << i + 1 << ": ";
        cin >> A[i];
    }

    // Ordenar ascendente mediante burbuja
    for (int i = 0; i < n - 1; i++) {
        for (int j = 0; j < n - i - 1; j++) {
            if (A[j] > A[j+1]) {
                aux   = A[j];
                A[j]  = A[j+1];
                A[j+1] = aux;
            }
        }
    }

    // Inicio del cronómetro: mide solo el bucle principal
    auto inicio = chrono::high_resolution_clock::now();

    do {
        total_generadas++;
        valida = true;

        for (int i = 0; i < n - 1; i++) {
            if (A[i] > 2 * A[i+1]) {
                valida = false;
                break;
            }
        }

        if (valida) {
            cout << "Permutación válida: ";
            for (int i = 0; i < n; i++)
                cout << A[i] << " ";
            cout << endl;
            total_validas++;
        }

    } while (next_permutation(A, A + n));

    auto fin = chrono::high_resolution_clock::now();
    chrono::duration<double> duracion = fin - inicio;

    cout << "Total permutaciones posibles:  " << factorial(n)        << endl;
    cout << "Total permutaciones generadas: " << total_generadas      << endl;
    cout << "Total permutaciones válidas:   " << total_validas        << endl;
    cout << "Tiempo de ejecución: "           << duracion.count()
         << " segundos" << endl;

    return 0;
}

// Calcula n! de forma iterativa
int factorial(int n) {
    int cont = 1;
    for (int i = 1; i < n; i++)
        cont = cont * (i + 1);
    return cont;
}

// Implementación propia de next_permutation (fines educativos)
bool next_permutation_test(int* inicio, int* fin) {
    int n = fin - inicio;
    if (n <= 1) return false;

    // Paso 1: buscar punto de quiebre de derecha a izquierda
    int i = n - 2;
    while (i >= 0 && inicio[i] >= inicio[i + 1]) i--;
    if (i == -1) return false;

    // Paso 2: buscar el sucesor de inicio[i]
    int j = n - 1;
    while (inicio[j] <= inicio[i]) j--;

    // Paso 3: intercambiar
    int temp  = inicio[i];
    inicio[i] = inicio[j];
    inicio[j] = temp;

    // Paso 4: invertir el sufijo desde i+1
    int izq = i + 1, der = n - 1;
    while (izq < der) {
        temp        = inicio[izq];
        inicio[izq] = inicio[der];
        inicio[der] = temp;
        izq++; der--;
    }
    return true;
}
\end{lstlisting}

% ══════════════════════════════════════════════════════════════════════════════
\section{Ejemplos de Entrada y Salida}
% ══════════════════════════════════════════════════════════════════════════════

Los siguientes ejemplos corresponden a ejecuciones reales del programa. Los tiempos
se midieron con \texttt{chrono::high\_resolution\_clock}.

\subsection{Ejemplo 1 --- Conjunto \{1, 3, 4\} (n=3)}

El conjunto del enunciado. Con tres elementos existen $3! = 6$ permutaciones.
El valor 4 supera el doble de 1, lo que descarta varias combinaciones.

\textbf{Salida real del programa:}
\begin{lstlisting}[style=pseudo]
=============================================
Permutación válida: 1 3 4
Permutación válida: 1 4 3
=============================================
Total permutaciones posibles:  6
Total permutaciones generadas: 6
Total permutaciones válidas:   2
Tiempo de ejecución: 0.0009147 segundos
\end{lstlisting}

\textbf{Verificación manual:}

\begin{center}
\begin{tabular}{ccccc}
\toprule
Permutación & Par (0,1) & Par (1,2) & Válida & Razón \\
\midrule
$[1, 3, 4]$ & $1 \leq 6$\;\checkmark & $3 \leq 8$\;\checkmark & Sí & Ambos pares cumplen \\
$[1, 4, 3]$ & $1 \leq 8$\;\checkmark & $4 \leq 6$\;\checkmark & Sí & Ambos pares cumplen \\
$[3, 1, 4]$ & $3 \leq 2$\;\texttimes & ---                    & No & $3 > 2 \cdot 1$     \\
$[3, 4, 1]$ & $3 \leq 8$\;\checkmark & $4 \leq 2$\;\texttimes & No & $4 > 2 \cdot 1$     \\
$[4, 1, 3]$ & $4 \leq 2$\;\texttimes & ---                    & No & $4 > 2 \cdot 1$     \\
$[4, 3, 1]$ & $4 \leq 6$\;\checkmark & $3 \leq 2$\;\texttimes & No & $3 > 2 \cdot 1$     \\
\bottomrule
\end{tabular}
\end{center}

\subsection{Ejemplo 2 --- Conjunto \{1, 2, 3\} (n=3)}

Con elementos consecutivos la restricción es menos estricta. Nótese que $[3, 2, 1]$
es válida porque $3 \leq 2 \cdot 2 = 4$ y $2 \leq 2 \cdot 1 = 2$: la igualdad cumple
la condición $\leq$.

\begin{lstlisting}[style=pseudo]
=============================================
Permutación válida: 1 2 3
Permutación válida: 1 3 2
Permutación válida: 2 1 3
Permutación válida: 3 2 1
=============================================
Total permutaciones posibles:  6
Total permutaciones generadas: 6
Total permutaciones válidas:   4
Tiempo de ejecución: 0.0014084 segundos
\end{lstlisting}

\subsection{Ejemplo 3 --- Conjunto \{1, 2, 3, 4\} (n=4)}

Con $n = 4$ el espacio de búsqueda crece a $4! = 24$ permutaciones y la proporción
de válidas baja al 33\,\%.

\begin{lstlisting}[style=pseudo]
=============================================
Permutación válida: 1 2 3 4
Permutación válida: 1 2 4 3
Permutación válida: 1 3 2 4
Permutación válida: 1 3 4 2
Permutación válida: 1 4 2 3
Permutación válida: 1 4 3 2
Permutación válida: 2 1 3 4
Permutación válida: 2 1 4 3
Permutación válida: 3 2 1 4
Permutación válida: 3 4 2 1
Permutación válida: 4 2 1 3
Permutación válida: 4 3 2 1
=============================================
Total permutaciones posibles:  24
Total permutaciones generadas: 24
Total permutaciones válidas:   12
Tiempo de ejecución: 0.0042051 segundos
\end{lstlisting}

\subsection{Ejemplo 4 --- Conjunto \{1, 5, 20\} (n=3) --- Elementos alejados}

Este caso ilustra cómo la distancia entre los valores del conjunto afecta el número
de válidas. Con elementos muy separados la restricción es muy estricta.

\begin{lstlisting}[style=pseudo]
=============================================
Permutación válida: 1 5 20
=============================================
Total permutaciones posibles:  6
Total permutaciones generadas: 6
Total permutaciones válidas:   1
Tiempo de ejecución: 0.00048 segundos
\end{lstlisting}

Comparando con $\{1, 2, 3\}$ que produjo 4 válidas, aquí solo se obtiene 1,
confirmando que la proporción de válidas depende tanto de $n$ como de los valores
específicos del conjunto.

% ══════════════════════════════════════════════════════════════════════════════
\section{Comparación de Tiempos de Ejecución}
% ══════════════════════════════════════════════════════════════════════════════

Se ejecutó el programa con el conjunto $\{1, 2, \ldots, n\}$ para distintos valores
de $n$. Los tiempos se obtuvieron con \texttt{chrono::high\_resolution\_clock}
midiendo únicamente el bucle principal.

\begin{center}
\begin{tabular}{rrrrr}
\toprule
$n$ & $n!$ & Válidas & Tiempo (s) & Factor \\
\midrule
6  & 720              & 144       & 0.0752    & ---      \\
8  & 40\,320          & 2\,880    & 2.0437    & $\times 27$x \\
10 & 3\,628\,800      & 86\,400   & 82.0305   & $\times 40$x \\
12 & 479\,001\,600    & 3\,628\,800 & 3\,962.63$^*$ & $\times 40$x \\
\bottomrule
\end{tabular}
\end{center}

{\small $^*$ Valor estimado por extrapolación a partir de los tiempos medidos para
$n = 6$, $n = 8$ y $n = 10$. La ejecución real de $n = 12$ supera los 45 minutos
y no fue completada, lo cual es consistente con la complejidad $O(n! \cdot n)$
del algoritmo.}

\bigskip
\textbf{Análisis de los resultados:}
\begin{itemize}
  \item \textbf{Ciclo $n=6$:} 720 permutaciones generadas en 0,075 segundos.
        El programa responde de forma instantánea. Se obtuvieron 144 permutaciones
        válidas (20\,\%).
  \item \textbf{Ciclo $n=8$:} 40\,320 permutaciones en 2,04 segundos. Tiempo
        perceptible pero práctico. Se obtuvieron 2\,880 válidas (7,1\,\%).
  \item \textbf{Ciclo $n=10$:} 3\,628\,800 permutaciones en 82 segundos. El tiempo
        ya es impracticable para uso interactivo. Se obtuvieron 86\,400 válidas (2,4\,\%).
  \item \textbf{Ciclo $n=12$:} con 479 millones de permutaciones el tiempo estimado
        supera los 45 minutos, confirmando que $n = 12$ es el límite práctico del algoritmo.
\end{itemize}

Se observa que el tiempo crece aproximadamente $\times 40$ cada dos niveles, lo que
es consistente con el factor $12 \times 11 = 132$ de crecimiento de $10!$ a $12!$.

% ══════════════════════════════════════════════════════════════════════════════
\section{Análisis de Complejidad}
% ══════════════════════════════════════════════════════════════════════════════

\subsection{Complejidad Temporal}

La complejidad temporal del algoritmo se descompone en dos partes:

\begin{itemize}
  \item El bucle principal itera exactamente $n!$ veces, una por cada permutación posible.
  \item Dentro de cada iteración, el bucle de validación recorre como máximo $n-1$
        pares consecutivos: $O(n)$.
\end{itemize}

Por lo tanto la complejidad temporal total es:

\begin{equation}
  T_{FB}(n) = O(n! \cdot n)
\end{equation}

Este es siempre el mismo costo sin importar el conjunto de entrada, ya que la fuerza
bruta siempre recorre las $n!$ permutaciones completas antes de descartar cualquiera.

\begin{center}
\begin{tabular}{rrrr}
\toprule
$n$ & $n!$ & $n! \cdot n$ & Operaciones aprox. \\
\midrule
8  & 40\,320       & 322\,560          & $3{,}2 \times 10^5$ \\
10 & 3\,628\,800   & 36\,288\,000      & $3{,}6 \times 10^7$ \\
12 & 479\,001\,600 & 5\,748\,019\,200  & $5{,}7 \times 10^9$ \\
\bottomrule
\end{tabular}
\end{center}

\subsection{Complejidad Espacial}

\begin{center}
\begin{tabular}{ll}
\toprule
Estructura & Espacio \\
\midrule
Arreglo $A$ de $n$ enteros & $O(n)$ \\
Variables auxiliares (contadores, aux, valida) & $O(1)$ \\
Permutaciones intermedias (no se almacenan) & $O(1)$ \\
\midrule
\textbf{Total} & $O(n)$ \\
\bottomrule
\end{tabular}
\end{center}

El algoritmo tiene complejidad espacial $O(n)$ porque en ningún momento almacena
todas las permutaciones en memoria: procesa cada una en el lugar y la descarta antes
de generar la siguiente.

% ══════════════════════════════════════════════════════════════════════════════
\section{Respuestas a las Preguntas Guía}
% ══════════════════════════════════════════════════════════════════════════════

\subsection{¿Qué ocurre con la proporción de permutaciones válidas cuando crece $n$?}

La proporción de permutaciones válidas \textbf{disminuye drásticamente} a medida que
$n$ crece. Esto ocurre porque la restricción $P[i] \leq 2 \cdot P[i+1]$ debe cumplirse
de forma simultánea en todos los $n-1$ pares consecutivos. Cuantos más pares haya que
satisfacer al mismo tiempo, menor es la probabilidad de que todos lo hagan.

Es importante notar además que el número exacto de válidas depende de los valores
específicos del conjunto, no sólo de $n$: conjuntos con elementos muy separados entre
sí generan menos válidas que conjuntos con elementos cercanos.

Usando los datos reales obtenidos con el conjunto $\{1, 2, \ldots, n\}$:

\begin{center}
\begin{tabular}{rrrrl}
\toprule
$n$ & Total ($n!$) & Válidas (medidas) & Proporción & Observación \\
\midrule
6  & 720              & 144       & 20,0\,\% & Medido experimentalmente \\
8  & 40\,320          & 2\,880    &  7,1\,\% & Medido experimentalmente \\
10 & 3\,628\,800      & 86\,400   &  2,4\,\% & Medido experimentalmente \\
12 & 479\,001\,600    & 3\,628\,800 & 0,5\,\% & Medido experimentalmente \\
\bottomrule
\end{tabular}
\end{center}

Se observa que la proporción disminuye de 20\,\% a 2,4\,\% al pasar de $n=6$ a
$n=10$, una reducción de casi un orden de magnitud en sólo 4 niveles. Mientras el
total crece factorialmente ($\times 5040$ de $n=6$ a $n=10$), las válidas crecen
mucho más lentamente ($\times 600$), haciendo que la proporción tienda a cero
conforme $n$ aumenta. Esto evidencia la ineficiencia creciente de la fuerza bruta.

\subsection{¿Dónde se podría aplicar poda con Backtracking?}

Con backtracking se construiría la permutación posición por posición, verificando la
restricción en cada paso antes de agregar el siguiente elemento. La poda ocurriría de
la siguiente manera:

\begin{itemize}
  \item En la posición 0 se prueba cada elemento disponible del conjunto.
  \item Antes de colocar un elemento en la posición $i$, se verifica si el elemento
        ya colocado en la posición $i-1$ cumple $A[i-1] \leq 2 \cdot A[i]$.
  \item Si no cumple $\Rightarrow$ se poda toda esa rama sin explorarla.
  \item Si cumple $\Rightarrow$ se continúa construyendo hacia la siguiente posición.
\end{itemize}

Por ejemplo con $\{1, 2, 3\}$, si en la posición 0 se coloca el 3:
\begin{itemize}
  \item Se prueba colocar 1 en posición 1: $3 \leq 2 \cdot 1 = 2$? \textbf{NO}
        $\rightarrow$ poda, no se exploran $[3, 1, 2]$ ni ninguna de sus variantes.
  \item Se prueba colocar 2 en posición 1: $3 \leq 2 \cdot 2 = 4$? \textbf{SÍ}
        $\rightarrow$ continúa explorando $[3, 2, \ldots]$.
\end{itemize}

La diferencia clave es que el backtracking nunca construye ni evalúa permutaciones
que ya se sabe de antemano que serán inválidas, mientras que la fuerza bruta siempre
las completa antes de descartarlas.

\subsection{¿Por qué \texttt{next\_permutation} genera exactamente $n!$ permutaciones?}

La función \texttt{next\_permutation} recorre las permutaciones en orden lexicográfico
estricto, una por una, sin repetir ni saltarse ninguna.

Al ordenar el arreglo de forma ascendente antes del bucle se garantizan dos condiciones:

\begin{enumerate}
  \item Se comienza desde la permutación mínima (la más pequeña en orden lexicográfico),
        por ejemplo $[1, 2, 3]$.
  \item Se termina en la permutación máxima (la más grande), por ejemplo $[3, 2, 1]$,
        momento en que \texttt{next\_permutation} retorna \texttt{false} y el bucle termina.
\end{enumerate}

Como entre la permutación mínima y la máxima de $n$ elementos distintos existen
exactamente $n!$ ordenamientos posibles, y \texttt{next\_permutation} los recorre todos
en orden sin repetir ninguno, se garantiza que genera exactamente $n!$ permutaciones.

Si el arreglo no estuviera ordenado inicialmente, el recorrido comenzaría desde la
mitad y sólo generaría las permutaciones mayores al estado inicial, perdiéndose las
anteriores y sin llegar a $n!$ en total.

% ══════════════════════════════════════════════════════════════════════════════
\section{Conclusiones}
% ══════════════════════════════════════════════════════════════════════════════

\begin{enumerate}
  \item El algoritmo de fuerza bruta genera correctamente las $n!$ permutaciones y
        filtra las que cumplen la restricción $P[i] \leq 2 \cdot P[i+1]$, produciendo
        resultados exactos en todos los casos probados.

  \item La complejidad temporal $O(n! \cdot n)$ hace que el algoritmo sea impracticable
        a partir de $n \approx 12$: el factorial crece mucho más rápido que cualquier
        función polinomial o exponencial de base fija, y para $n = 13$ se superan los
        $6 \times 10^9$ ciclos de trabajo.

  \item La proporción de permutaciones válidas disminuye drásticamente con $n$: usando
        $\{1, 2, \ldots, n\}$, las válidas crecen como potencias de 2 mientras el total
        crece factorialmente. Para $n = 7$ la proporción ya es menor al 2\,\%, lo que
        evidencia la ineficiencia creciente de la fuerza bruta.

  \item El número exacto de válidas depende tanto de $n$ como de los valores específicos
        del conjunto: elementos muy separados entre sí reducen las válidas más
        drásticamente que elementos cercanos.

  \item La función \texttt{next\_permutation} garantiza exactamente $n!$ iteraciones
        siempre que el arreglo esté ordenado ascendentemente al inicio, condición
        necesaria para el correcto funcionamiento del algoritmo.

  \item Un enfoque de backtracking podría mejorar significativamente la eficiencia al
        podar ramas inválidas durante la construcción de cada permutación, evitando
        explorar subárboles completos de soluciones que ya se sabe que serán inválidas.
\end{enumerate}

% ══════════════════════════════════════════════════════════════════════════════
\begin{thebibliography}{9}
% ══════════════════════════════════════════════════════════════════════════════

\bibitem{cormen}
T.~H. Cormen, C.~E. Leiserson, R.~L. Rivest, C.~Stein,
\textit{Introduction to Algorithms}, 4th~ed. MIT Press, 2022.

\bibitem{levitin}
A.~Levitin,
\textit{Introduction to the Design and Analysis of Algorithms}, 3rd~ed.
Pearson, 2012.

\bibitem{skiena}
S.~S. Skiena,
\textit{The Algorithm Design Manual}, 2nd~ed. Springer, 2008.

\bibitem{cppreference}
Cppreference, \textit{std::next\_permutation}, 2024.
Disponible en: \url{https://en.cppreference.com/w/cpp/algorithm/next_permutation}

\bibitem{alvarez}
C.~A. Álvarez Henao,
\textit{Material de clase --- Análisis y Diseño de Algoritmos}.
Universidad EAFIT, 2026.

\end{thebibliography}

\end{document}